% Nouveau EGA V (cote 25): archive cote 25, site page 128 (PDF page 129).
% typescript. Draft compared: none.
% AI-transcribed by Claude Opus 5.5 (claude-opus-5-5, adaptive thinking, effort high), Message Batches API, batch ega25b (msgbatch_01A7n5YuXubUPNK1zsepi8RG), 2026-11-13; not yet checked by eye.
\scanpage{25/128}
\mspage{\hand{31}}

par une section inversible de $\underline{O}_S$ , donc $\hand{\phi}$ et $\hand{\psi}$ définissent le même point
de $\overset{\hand{\vee}}{P}$ à valeurs dans S. Comme tout point de $\overset{\hand{\vee}}{P}$ à valeurs dans S est défini,
localement sur S, par une section $\hand{\phi}$ de $\underline{E}_S$ qui ne s'annule pas (cf Chap II)
a) s'ensuit.   Pour prouver b), on note le

\uline{Lemme}   Soit D un syst lin de diviseurs sur X non vide loc de type fini
sur k alg clos, paramétré par $\overset{\hand{\vee}}{P}$(E), et considérons le morphisme correspon-%
dant  f: X-Z $\hand{\longrightarrow}$ P , où Z est l'ensemble des points fixes. Alors  si
r = rang$_k$E \hand{$>0$,} il existe  r+1 point\hand{s}$x_i$$(1\leq i\leq r+1)$ de  X(k)-Z(k)  tels que les
f(x$_i$) forment une  ``base projective'' de P, i.e. tels que  pour toute
partie J de [1,r+1] ayant \hand{r}  éléments,  les  f(x$_i$),i\hand{$\in J$},  ne sont pas
contenus dans un hyperplan de P.

On peut évidemment supposer Z=\hand{$\emptyset$} ,   Comme par \unsure{3|5}.1. f(X) n'est contenu dans
aucun hyperplan de P, on en conclut d'abord l'existence de r points
$x_i$$(1\leq i\leq r)$ de X \unsure{teks|tels} que les f($x_i$) soient projectivement indépendants dans P, i.e. soient définies
par des formes linéairement indépendantes sur E.   Reste à prouver qu'il
existe un  $x_{r+1} = x$   dans X(k) tel que  f(x) ne soit dans aucun des r
\unsure{hyperpbans|hyperplans} $H_i$ définis par les  systèmes de r-1 parmi les f(x$_i$). Or dans
le cas contraire, compte tenu du sorite Jacobson,  on aurait
f(X) $\hand{\subset\bigcup}_{\hand{i}}$ H$_i$ , donc  si X$_0$  est une composante
irréductible de X, f(X$_0$) serait contenu dans un des H$_i$ , ce qui contredit
\unsure{3|5}.1. - Ceci posé,  pour prouver b), on peut  évidemment supposer Z=\hand{$\emptyset$} ,
et utilisant \unsure{3|5}.1.  on est ramené à prouver qu'un automorphismeu de $\overset{\hand{\vee}}{P}$ est
déterminé quand on connait le composéde son contragrédient $\overset{\hand{\vee}}{u}$ dans P avec f: X$\hand{\longrightarrow}$ P , - et que l'asser-%
tion analogue est vraie après tout changement de base S$\hand{\rightarrow}$Spec(k)
pour un  automorphisme u de $\overset{\hand{\vee}}{P}_S$  .  Or cela résulte aussitôt  du lemme

\mspage{\hand{128}}
